\section{A priori estimates for the deformation}
\label{n3:apriori:section}

We bound the height, boundary slopes, and conformal variable before
using classical regularity. Fix the prepared exterior of
Proposition~\ref{n3:domains:prepared-domain}, its thresholds and collars,
and $\epsilon>0$. We use the variables and four stages of
Proposition~\ref{n3:deformation:homotopy}. A solution is smooth on a finite
truncation and satisfies $t\ge-\epsilon$, so the estimates include
putative floor-contact solutions. No Hessian bound or upper bound for
$Z$ is assumed.

A quantity denoted by $\Pi_N$
is bounded by $C(1+N)^A$, with $C,A$ independent of the truncation,
the homotopy parameter, and $N$.  It can increase at successive uses.
The fixed data, collars, and $\epsilon$ can enter $C,A$.  A constant
denoted by $C(N)$ can depend arbitrarily on these fixed data and on
$N,\ell,\tau$; it is still independent of the truncation and the
homotopy parameter.  Constants explicitly described as fixed do not
depend on $N$.  Recall that
\begin{equation}
 \ell=e^{-\epsilon N},\qquad \tau=\ell^{3/2},\qquad
 \ell\le l\le e,\qquad 0\le p\le N,
 \qquad \frac{4d}{4+N}\le\chi\le d.
 \label{n3:apriori:parameter-comparisons}
\end{equation}
In particular $\Pi_N\tau/\ell\to0$. This small factor will absorb
only terms whose constants are polynomial.

\subsection{Sublevels and exclusion of boundary layers}

The limiting height will be only lower semicontinuous.  We first show
how an expansion bound on its lower tests controls every exterior
support of a closed sublevel, including a plateau.  This local step
will also apply to the four-dimensional constructions.

\begin{lemma}[Exterior supports of a closed sublevel]
\label{n3:apriori:sublevel-support}
Let $U$ be an open subset of a smooth Riemannian manifold of dimension
$d\in\{3,4\}$, let $Q$ be a smooth symmetric tensor on $U$, and let
$u:U\to\R$ be finite and lower semicontinuous.  Fix $k<k'$ and
$q_0\in\R$.  A smooth lower test of $u$ at $x$ is a smooth function
$\phi$ with $\phi(x)=u(x)$ and $\phi\le u$ near $x$.
Suppose every such test with $u(x)<k'$ and $\dd\phi(x)\ne0$ satisfies
\[
 \mathcal E_Q[\phi](x)\le q_0,
\]
where $\mathcal E_Q$ is the expansion in
Equation~\eqref{n3:domains:test-expansion} of the regular level through
$x$, oriented toward increasing $\phi$.  Equivalently, subtract
$\phi(x)$ before applying that formula.
Then every smooth exterior support of the relatively closed set
$D_k=\{x\in U:u(x)\le k\}$, at a frontier point in $U$, has
expansion at most $q_0$.
\end{lemma}

\begin{proof}
For integers $j\ge1$ define
\begin{equation}\label{n3:apriori:sigmoid}
 S_j(s)=\bigl(1+\exp[-j^2(s-k-j^{-1})]\bigr)^{-1},
 \qquad v_j=S_j\circ u.
\end{equation}
These functions are lower semicontinuous and take values in $(0,1)$.
Their lower relaxed limit $\liminf_{j\to\infty,\,y\to x}v_j(y)$
is the function $I_k(x)$ equal to zero on $D_k$ and one on its
complement.  Indeed, at $x\in D_k$ the constant
sequence gives $v_j(x)\le S_j(k)\to0$, while $v_j\ge0$ everywhere.
If $u(x)>k$, lower semicontinuity gives a neighborhood on which
$u\ge k+\delta$ for some $\delta>0$, and $v_j\to1$ uniformly there.

Let $\phi$ define an exterior support of $D_k$ at $x$:
$\phi(x)=0$, $\dd\phi(x)\ne0$, and $\phi\le0$ on $D_k$ near $x$.
Shrink to a small normal-coordinate ball compactly contained in $U$ on which
$|\phi|<1/2$.  Then $\phi$ is a lower test of $I_k$ at $x$.
Subtract a small positive multiple of the fourth power of the distance
from $x$ to make this contact strict without changing its two-jet;
write the resulting function as $\phi_0$.
Minimize $v_j-\phi_0$ on the closed ball.  The preceding relaxed-limit
description, the strict contact, and $v_j(x)\to0$ show that the
minimizers $x_j$ lie in its interior for large $j$, satisfy $x_j\to x$,
and have contact values $v_j(x_j)\to0$.  Thus $\phi_0+a_j$ is a
lower test of $v_j$ at $x_j$, where $a_j\to0$.

For each fixed large $j$, restrict to a neighborhood of $x_j$ on
which $\phi_0+a_j$ takes values in $(0,1)$.  The increasing function
$S_j^{-1}$ makes $S_j^{-1}(\phi_0+a_j)$ a smooth lower test of $u$.
Its contact height is
\[
 u(x_j)=k+j^{-1}
       +j^{-2}\log\frac{v_j(x_j)}{1-v_j(x_j)}
       \le k+j^{-1}<k'
\]
for all sufficiently large $j$.  Its gradient is nonzero, since
$\dd\phi_0(x_j)\to\dd\phi(x)\ne0$.
An increasing smooth composition preserves the oriented level expansion:
it preserves the unit normal, and for $\theta'>0$,
\[
 \Hess(\theta\circ\phi_0)
  =\theta'\Hess\phi_0
       +\theta''\dd\phi_0\otimes\dd\phi_0,
 \qquad
 \frac{\tr_{(\dd\phi_0)^\perp}\Hess(\theta\circ\phi_0)}
 {|\dd(\theta\circ\phi_0)|}
 =\frac{\tr_{(\dd\phi_0)^\perp}\Hess\phi_0}{|\dd\phi_0|}.
\]
The hypothesis therefore gives
$\mathcal E_Q[\phi_0+a_j](x_j)\le q_0$.
Letting $j\to\infty$ yields $\mathcal E_Q[\phi](x)\le q_0$,
because $\phi_0$ and $\phi$ have the same two-jet at $x$.
\end{proof}

For notation in this section, let $\mathcal B_c$ be the full closed
maximal black region at threshold $c$, and let $\mathcal W_c$ be the
corresponding white region with $\mathcal B_{c_b}$ held fixed.  These
are the families used in Proposition~\ref{n3:domains:prepared-domain}.

\begin{lemma}[Uniform separation from the wrong height]
\label{n3:apriori:sublevels}
In the first two homotopy stages, the following holds.  If
$Q\subset\bar\Omega$ is compact and disjoint from $\mathcal B_{c_b}$,
there are $\eta_Q>0$ and $N_Q$ such that
\begin{equation}
 h\ge b_-+\eta_Q\quad\hbox{on }Q
 \label{n3:apriori:lower-height}
\end{equation}
for every $N\ge N_Q$, every sufficiently large allowed truncation,
and every admissible or floor-contact solution in those stages.
If $Q$ is disjoint from $\mathcal W_{c_w}$, the analogous conclusion is
$h\le b_+-\eta_Q$.  The compact sets are relative to the full closure:
they may contain faces of the opposite color.
\end{lemma}

\begin{proof}
We prove the lower assertion.  It is enough to consider an arbitrary
sequence $N_i\to\infty$, $R_i\to\infty$ and solutions in either of the
first two stages.  The elementary height and end estimates in
Lemmas~\ref{n3:deformation:height} and~\ref{n3:deformation:end-height} give
\begin{equation}
 |h_i|\le C_0,\qquad
 |h_i|\le C(r^{-\gamma}-R_i^{-\gamma})\quad\hbox{on the end},
 \qquad 3/4<\gamma<1.
 \label{n3:apriori:elementary-height}
\end{equation}
Near each white face extend $h_i$ to the other side by the constant
$b_+$.  Extend the fixed geometric coefficients smoothly there.
The extended functions are continuous.  Define the lower relaxed limit
\[
 \underline h(x)=\lim_{j\to\infty}
       \inf\{h_i(y):i\ge j,\ \operatorname{dist}(x,y)<1/j\}.
\]
It is lower semicontinuous, bounded, and satisfies the limiting end
bound.  No equicontinuity of $h_i$ is asserted.

Choose $b_-<k'<\min\{0,b_-+\eta\}$, where $\eta$ is the low-height
sign range of the added term $D_0$.  Suppose that a smooth function
$\phi$ touches $\underline h$ from below at $x$, with
$\underline h(x)\le k'$ and $d\phi(x)\ne0$.  Exclude for the moment
the black faces.  Subtracting a small fourth-order localization term
from $\phi$ makes the contact strict without changing its two-jet.
Minimization on a small closed ball then produces, after passing to a
subsequence, points $x_i\to x$ and constants $c_i\to0$ such that
$\phi+c_i$ touches $h_i$ from below
at $x_i$, and $h_i(x_i)\to\underline h(x)$.  The boundary of this
ball cannot contain $x_i$ by strictness.  Nor can $x_i$ lie on a white
face or on the extended side: their values there are $b_+>k'$.
Thus every sufficiently large $i$ uses the interior trace equation.

At these contacts,
\[
 \nabla f_i=\tau_i^{-1}\nabla\phi,
 \qquad
 d_i\le\frac{\tau_i^2}{\ell_i^2|d\phi(x_i)|^2}
          \longrightarrow0,
 \qquad a_i\longrightarrow1,
 \qquad \chi_i\longrightarrow0.
\]
Here $a_i=l_i|\nabla f_i|/\sqrt{D_i}$ is the deformation variable.
Positivity of $A_{\chi_i}$ and the Hessian
inequality at a lower contact imply
\[
 \tr_{A_{\chi_i}}K+
 \frac{l_i}{\tau_i\sqrt{D_i}}
                 \tr_{A_{\chi_i}}\Hess\phi
 \le F_i(x_i).
\]
In the second stage $D_0\le0$ at these values; in the first stage it
is absent.  Consequently $\limsup F_i(x_i)\le k'+C_b$.
Passing to the limit gives
\begin{equation}
 \frac{\tr_{(d\phi)^\perp}\Hess\phi}{|d\phi|}
           +\tr_{(d\phi)^\perp}K\le k'+C_b.
 \label{n3:apriori:relaxed-test}
\end{equation}
This is the expansion of the level surface of $\phi$, oriented toward
increasing $\phi$.  The only Hessian used in this passage is the
fixed test Hessian.  Removing the localization proves the assertion
for non-strict contacts as well.

Fix $b_-<k<k'$ and put $E_k=\{\underline h\le k\}$.
Apply Lemma~\ref{n3:apriori:sublevel-support} with
$u=\underline h$, tensor $Q=K$, and $q_0=k'+C_b$.
The function is finite and lower semicontinuous, and
Equation~\eqref{n3:apriori:relaxed-test} supplies the required bound
for every nonzero-gradient lower test below $k'$.  We apply the lemma
locally away from black faces, using the constant extensions to obtain
open neighborhoods across white faces.  Every exterior support of
$E_k$ there consequently has expansion at most $k'+C_b$, with no
regularity assumption on the sublevel.

This argument also rules out a hidden layer at a white face.
At every point strictly on the added side, the relaxed limit equals
$b_+>k$.  Thus $E_k$ lies locally in the original closed exterior,
and the reversed white face is an exterior support of $E_k$ at any
contact there, with expansion
\[
 -H_B+P_B=-c_w>0.
\]
For $k'$ sufficiently close to $b_-$, this is strictly greater than
$k'+C_b$.  The high extension ensured that all preceding low-value
tests came from the interior equation, so the support contradiction
applies at the face itself.  Hence $E_k$ avoids every white face.
The end estimate makes $E_k$ compact.

Adjoin the entire closed region $\mathcal B_{c_b}$.
At a point on its smooth black frontier,
every exterior support of the union also contains the black region
locally.  The tangency comparison therefore bounds its expansion by
$c_b<k'+C_b$.  At all other frontier points the preceding sublevel
argument applies.  The resulting compact closed set
contains the required original inner barrier with a collar, because
the black region already does.  It avoids the outer barriers.
Lemma~\ref{n3:domains:weak-support} places it in a smooth black trapped
region at any threshold $c''>k'+C_b$ sufficiently close to $c_b$.
It is consequently contained in $\mathcal B_{c''}$.

Finally let $Q$ be as in the lemma.  Hausdorff right continuity at
$c_b$ gives a $\delta>0$ such that
$\mathcal B_{c_b+\delta}\cap Q=\varnothing$; decrease $\delta$ so
that the enlarged threshold is still negative and all barriers
remain strict.  Choose $k,k',c''$ with
\[
 b_-<k<k',\qquad k'+C_b<c''<c_b+\delta.
\]
The construction implies $E_k\cap Q=\varnothing$ for every sequence
under consideration.  If the uniform estimate $h>k$ on $Q$ failed,
choose violating solutions with $N_i\to\infty$, $R_i\to\infty$ and
$x_i\in Q$, $h_i(x_i)\le k$.  A subsequence has $x_i\to x\in Q$
and $\underline h(x)\le k$, a contradiction.  This gives
\eqref{n3:apriori:lower-height}, after decreasing its positive margin.

For the upper assertion apply the entire argument to
$(-h,-K,-C)$, interchanging colors.  The reversed black faces have
strictly positive white expansion $-c_b$.  Their exclusion and
compactness give positive distance from the black region before
applying the white version of Lemma~\ref{n3:domains:weak-support}.
Adjoining the full white region is therefore legitimate in the
fixed black complement.  Empty white regions and missing face types
cause no change: only an actually nonempty sublevel would supply a
trapped seed and hence a contradiction.
\end{proof}

For the truncation quantifier, start with any $R_{\min}(N)\to\infty$
that satisfies the outer estimates of
Section~\ref{n3:deformation:section}, and enlarge it when necessary.
Failure of any asserted estimate for arbitrarily large $N$ and
$R\ge R_{\min}(N)$ supplies exactly the sequence excluded above.
Only finitely many compact sets will be needed below, so their
thresholds can be chosen simultaneously.

The weak-set argument has kept low heights near black faces and high
heights near white faces. We now sharpen this separation into a
nonzero signed boundary slope. That slope makes the otherwise
uncontrolled inner flux small enough for the mass estimate.

\subsection{Signed collar barriers}

\begin{lemma}[Boundary slopes]
\label{n3:apriori:collars}
There are fixed constants $c,C_*>0$ such that, for sufficiently large
$N$ and the allowed truncations, the first two stages satisfy
\begin{equation}
 \begin{array}{ll}
 c/\tau\le\partial_\nu f\le C_*/\tau&\text{on black faces},\\[2pt]
 c/\tau\le-\partial_\nu f\le C_*/\tau&\text{on white faces}.
 \end{array}
 \label{n3:apriori:signed-slopes}
\end{equation}
During the third stage one still has
$|\partial_\nu f|\le C_*/\tau$.  On every fixed compact region needed
for the final scalar-curvature estimate, the first two stages satisfy
$b_-\le h\le b_+$.
\end{lemma}

\begin{proof}
Write $s$ for the smooth leaf parameter in an outward collar and
$\nu_s=\nabla s/|ds|$.  All geometric coefficients of this fixed collar
are bounded, with $|ds|$ bounded above and below.  At a comparison
contact with $h_0=b_-+\psi(s)$, $\psi'>0$, the shared gradient gives
\begin{equation}
 d\le C\frac{\tau^2}{\ell^2(\psi')^2},\qquad
 1-a\le d,\qquad N\chi\ge c_1d\quad(N\ge4).
 \label{n3:apriori:contact-degeneracy}
\end{equation}
If the slopes are in a fixed positive interval, the first bound is
$d\le C\ell$.  None of these comparisons uses an upper bound for $Z$.

The trace operator on this test is
\begin{align}
 &\tr_{A_\chi}\left(K+
       \frac{l}{\tau\sqrt D}\Hess h_0\right)\notag\\
 &\quad=P_s+aH_s+\chi K(\nu_s,\nu_s)
       +a\chi\left(
           \frac{\Hess s(\nu_s,\nu_s)}{|ds|}
                      +|ds|\frac{\psi''}{\psi'}\right),
 \label{n3:apriori:collar-trace}
\end{align}
where $P_s$ is the tangential trace of $K$ on the leaf.  Thus it is
the leaf expansion $H_s+P_s$, with an error of absolute value at most
$C d$, plus the displayed logarithmic-slope term.  The same formula
for a negative slope has $a$ replaced by $-a$ in the Hessian terms.
This is an evaluation at the contact variables; it does not
differentiate the unknown $t$.

On a black collar $H_s+P_s\ge c_b$ and $C=C_b-k_Bs$.
Choose small positive slopes so that $\psi(0)=0$ and
$\psi(s)\le k_Bs/2$.  At the outer end of a sufficiently short fixed
collar, Lemma~\ref{n3:apriori:sublevels} lets this test lie below the
solution.  In the compatible direction the add-on $D_0$ vanishes.
The residual of the trace equation on the test is therefore at least
\begin{equation}
 k_Bs/2-Cd+a\chi|ds|\frac{\psi''}{\psi'}.
 \label{n3:apriori:lower-collar-residual}
\end{equation}
Choose a smooth nonnegative $\beta_N(s)$ equal to $AN$ for
$0\le s\le1/N$, supported in $0\le s\le2/N$, and satisfying
$\int\beta_N\le2A$.  Define
\[
 \psi'(s)=m\exp\left(\int_0^s\beta_N(q)\dd q\right).
\]
Choose the fixed $A$ large enough that the final term of
\eqref{n3:apriori:lower-collar-residual} dominates $Cd$ on $s\le1/N$,
using \eqref{n3:apriori:contact-degeneracy} and $a\ge1/2$ there.
On $s\ge1/N$, the first term dominates $Cd\le C\ell$ for large $N$;
the taper contributes with the good sign.  Since the slope multiplier
is at most $e^{2A}$, a fixed sufficiently small $m>0$ enforces all
the previous small-slope requirements.  This constructs a strict
lower barrier with slopes bounded above and below independently of
$N$.

An upper barrier from $b_-$ is constructed in the same collar with
\[
 \psi'(s)=M\exp\left(-\int_0^s\beta_N(q)\dd q\right).
\]
The smooth leaf expansion differs from its boundary value by $O(s)$.
Choose the fixed $M$ so large that
$\psi(s)\ge C_2s$ dominates this variation, the offset variation,
and the height bound at the outer collar end.  The logarithmic-slope
term is now negative.  On $s\le1/N$ it dominates $Cd$; on the rest
of the collar the negative margin proportional to $s$ does so.
The upper barrier is strict at every possible positive contact.
Its slopes lie in $[Me^{-2A},M]$, an interval independent of $N$.

To check the comparison signs, at a negative interior minimum of
$h-h_0$ the Hessian of $h$ dominates that of $h_0$, the gradients
and $Z$ agree, and hence so do $t,l,A_\chi$.  The right side is
strictly increasing in the height.  Thus a strict positive residual
for a lower barrier contradicts the equation.  At a positive maximum
the inequalities reverse, proving the upper-barrier assertion.
Taking the one-sided derivative at $s=0$ proves the black part of
\eqref{n3:apriori:signed-slopes}.  Replacing $(h,K,C)$ by
$(-h,-K,-C)$ proves the white part.

In the third stage the boundary height is the assigned interpolated
constant, say $b_\zeta$.  Construct tests
$b_\zeta+\psi(s)$ and $b_\zeta-\psi(s)$ with large positive
$\psi'$ and $(\log\psi')'=-\beta_N$.  At $s=0$, the two directional
inequalities of Proposition~\ref{n3:deformation:homotopy} say precisely
that the positive-slope test has nonpositive limiting residual and
the negative-slope test nonnegative limiting residual.  At every
comparison contact, the fixed positive barrier slope and
Equation~\eqref{n3:apriori:contact-degeneracy} give $d\le C\ell$, so
$a\ge1/2$ for large $N$.  On $1/2\le|w|\le1$, the first
$w$-derivatives of $A_\chi$ in
Equation~\eqref{n3:deformation:axial-matrices} are bounded independently
of $0\le p\le N$.  The other collar terms have fixed smooth
coefficients on the bounded height range.  Thus, at fixed boundary
height the coefficient and geometric errors are $O(s+d+\chi)$,
since $|w-(\pm\nu_s)|=1-a\le d$ and the radial comparison segment stays
in this range.  Changing the test height supplies, by height
monotonicity, a residual of magnitude at least $\psi(s)$ with the
required sign.  Large fixed slopes dominate the $O(s)$ error and
the outer-end height bound.  The logarithmic-slope term has the
required sign for both tests, as follows from
\eqref{n3:apriori:collar-trace}; it dominates the $O(d)$ error on the
initial $1/N$ interval.  The remaining margin dominates it beyond
that interval.  This proves the two-sided upper derivative bound,
without using sublevel exclusion in the third stage.

Finally, on each own-color collar the lower barrier just constructed
gives the required one-sided height bound.  On the rest of a fixed
compact region use Lemma~\ref{n3:apriori:sublevels}; the opposite height
bound on an own-color collar follows from that lemma as well, since
the closed collar is disjoint from the other region.  A finite cover
proves $b_-\le h\le b_+$ on the compact regions in the statement.
\end{proof}

\subsection{A trace estimate with polynomial constants}

The collar trace estimate must have polynomial constants: it will
control the boundary flux both here and in the mass estimate.
We therefore prove it before comparing with the ordinary product graph.

\begin{lemma}[Weighted and unweighted collar traces]
\label{n3:apriori:trace}
In the first two stages let $\mathcal C$ be a fixed union of disjoint
collars of the inner faces, shortened if necessary.  For every
nonnegative smooth function $\varphi$ supported in those collars,
\begin{align}
 \int_B L\varphi\dd A_g
 &\le\Pi_N\int_{\mathcal C}u
       \bigl[(1+\sqrt{\mathcal T})\varphi
                          +|d\varphi|_{A_\chi}\bigr]\dd V_g,
 \label{n3:apriori:weighted-trace}\\
 \int_B\varphi\dd A_g
 &\le\Pi_N\int_{\mathcal C}\sqrt D
       \bigl[(1+\sqrt{\mathcal T})\varphi
                          +|d\varphi|_{A_\chi}\bigr]\dd V_g.
 \label{n3:apriori:unweighted-trace}
\end{align}
The estimates also hold without a support condition on $\varphi$
if the right sides include a fixed collar cutoff equal to one on
$B$, with its derivative included in the coefficient of $\varphi$.
\end{lemma}

\begin{proof}
For $\sigma>0$, write $e_f=\nabla f/\sigma$, the axis denoted by
$e$ in Section~\ref{n3:deformation:section}. Set
\[
 W=l\bar\nabla f=\frac{l\nabla f}{D}=\sqrt d\,a e_f.
\]
Its norm in $\bar g$ is $a\le1$.  For any covector $\xi$,
\begin{equation}
 |\xi(W)|\le |\xi|_{A_d}
                  \le\sqrt{1+N/4}\,|\xi|_{A_\chi}.
 \label{n3:apriori:W-gradient}
\end{equation}
Direct differentiation, using $uW=Lw$, gives
\begin{align}
 \frac{\Div(uW)}u
 &=\sqrt d\left[
       \tr_{A_d}H^f+(dp+p_L)a\partial_{e_f}t\right],
 \label{n3:apriori:W-divergence}\\
 \frac{\Div(\sqrt D\,W)}{\sqrt D}
 &=\sqrt d\left[
       \tr_{A_d}H^f+dp\,a\partial_{e_f}t\right].
 \label{n3:apriori:W-divergence-unweighted}
\end{align}
For example,
$\Div w=\tr_{A_d}H^f+dp\,a\partial_{e_f}t$,
while differentiating $L$ adds $p_La\partial_{e_f}t$ to the first
formula.  Coercivity in Lemma~\ref{n3:deformation:coercivity}
bounds both right sides by $\Pi_N(1+\sqrt{\mathcal T})$.
More explicitly, writing $q=(K+H^f)_{e_fe_f}$, the normal Hessian
term has coefficient $d^{3/2}$, so it is bounded by $\sqrt d|q|+C$.
The other Hessian terms are transverse traces, and the gradient
term is at most $(2N+2)\sqrt d|\partial_{e_f}t|$.
All are controlled by the stated coercivity estimate.  This also
establishes the formulas and bounds at zero gradient by their smooth
direction-free extensions.

Let $\varepsilon_B=1$ on a black collar and $-1$ on a white collar.
The signed slopes give $w_\nu=\varepsilon_B a$ on $B$ and
$a\ge1/2$ for large $N$.  Thus
\[
 uW_\nu=\varepsilon_B La,
 \qquad \sqrt D\,W_\nu=\varepsilon_B a.
\]
The outward integration normal is $-\nu$.
Apply the divergence theorem to $-\varepsilon_B uW\varphi$, with
a fixed collar cutoff.  Its boundary flux is $La\varphi$.
Equations~\eqref{n3:apriori:W-gradient} and
\eqref{n3:apriori:W-divergence}, and the fixed cutoff derivative, give
\eqref{n3:apriori:weighted-trace}.  Replace $u$ by $\sqrt D$ and use
\eqref{n3:apriori:W-divergence-unweighted} for
\eqref{n3:apriori:unweighted-trace}.  No factor $\ell^{-1}$ was used.
\end{proof}

\subsection{High-level energy and the first upper bound}

The signed slopes and polynomial trace estimate are now available.
We use them before allowing any constants that depend arbitrarily on
$N$. The first stage contains the drift in the divergence equation. We first
bound this stage uniformly for $\lambda\in[0,1]$.
There is a number $M_0=M_0(N)>0$ such that
\begin{equation}
 \Xi\ge\delta_0\mathcal T+\rho+
                      \tfrac12\delta_0\tau a\sigma
          \quad\hbox{on }\{Z>M_0\}.
 \label{n3:apriori:high-source}
\end{equation}
Indeed, outside the support of $m_0$ this is immediate.  On its
support, $-\epsilon\le t\le-\epsilon+1/N$, so $l$ is comparable
to $\ell$.  The identity $D=e^{8(Z-t)}$ shows that $a\sigma$ tends
uniformly to infinity as $Z\to\infty$ in that band.  The penalty
$\mathcal P$ is bounded on the whole exterior by $\Pi_N$.
It is therefore absorbed by $\delta_0\tau a\sigma/2$ above a fixed
$M_0(N)$.  This choice is independent of $R$ and $\lambda$.

Choose a smooth nondecreasing $k_0$ equal to zero on $Z\le M_0$,
equal to one on $Z\ge M_0+1$, and with bounded derivative.  Test
$\Div V_\lambda=u\Xi$ by $k_0(Z)$.  Its outer boundary term
vanishes since $Z=0$.  The drift cross term is bounded by
\[
 u k_0'\left|\langle dZ,K(w,\cdot)\rangle_{A_\chi}\right|
 \le 2u k_0'|dZ|_{A_\chi}^2+C u k_0'|K|^2.
\]
The last term has bounded integral: it is supported on a fixed
compact set and on $M_0\le Z\le M_0+1$, where the floor and
$D=e^{8(Z-t)}$ bound $u$ and $\sigma$ by $C(N)$.

At a black face $F-P_B=H$, $w_\nu=a$; at a white face
$F-P_B=-H$, $w_\nu=-a$.  In both cases
\[
 -H+w_\nu(F-P_B)=(a-1)H.
\]
The additional omitted-drift term in the homotopy boundary condition
has absolute value at most $C\chi\le Cd$, since
$A_\chi K(w,\cdot)$ has normal component
$\chi w_\nu K(\nu,\nu)$.  The remaining term is $-N_Bd$, and the
boundary multiplier lies between zero and one.  Hence
\begin{equation}
 |(V_\lambda)_\nu|\le Cud=CL\sqrt d
                       \le C\frac\tau\ell L\quad\hbox{on }B,
 \label{n3:apriori:small-boundary-flux}
\end{equation}
where the last inequality uses the lower slope in
\eqref{n3:apriori:signed-slopes}.  Integration now gives
\begin{align}
 &\int_{\Omega_R}u k_0
       \left(\delta_0\mathcal T+\rho+
                      \tfrac12\delta_0\tau a\sigma\right)\dd V_g
       +2\int_{\Omega_R}u k_0'|dZ|_{A_\chi}^2\dd V_g\notag\\
 &\hspace{18mm}\le C(N)+C\frac\tau\ell
                                  \int_B Lk_0\dd A_g.
 \label{n3:apriori:high-energy-before-trace}
\end{align}

Apply \eqref{n3:apriori:weighted-trace} with the fixed collar cutoff.
The boundary cost is at most
\[
 \Pi_N\frac\tau\ell\int_{\mathcal C}u
       \left[k_0(1+\sqrt{\mathcal T})
                              +k_0'|dZ|_{A_\chi}\right]\dd V_g.
\]
Let $\rho_{\mathcal C}>0$ be the minimum of $\rho$ on the closed
collars.  Pointwise,
$1+\sqrt{\mathcal T}\le C(\rho+\delta_0\mathcal T)$ there, with
$C$ fixed.  Since $\Pi_N\tau/\ell\to0$, the first summand is
absorbed in the first integral of
\eqref{n3:apriori:high-energy-before-trace}.  Young's inequality bounds
the derivative summand by
\[
 u k_0'|dZ|_{A_\chi}^2+
          C(\Pi_N\tau/\ell)^2u k_0'.
\]
The second term again has bounded compact integral because it lies
in the transition strip of $k_0$.  We obtain
\begin{equation}
 \int u k_0(\mathcal T+\rho+\tau a\sigma)\dd V_g
       +\int u k_0'|dZ|_{A_\chi}^2\dd V_g\le C(N).
 \label{n3:apriori:high-energy}
\end{equation}
The smallness requirement on $N$ in this absorption is independent
of $M_0$; only its resulting right-hand constant depends on $M_0$.

The polynomial boundary absorption is complete. We may now use
arbitrary fixed-$N$ comparison constants. Write
$\Gamma$ for the ordinary product graph of $f$ in $g+(\dd b_0)^2$,
and $\dd\Gamma=\sqrt{1+\sigma^2}\dd V_g$.  The measures
$\dd\Gamma$ and $\sqrt D\dd V_g$ are comparable by constants
depending only on $\ell$.  On every fixed compact set $Q$,
\begin{equation}
 \int_{\Gamma\cap\pi^{-1}(Q)}e^{2Z}\dd\Gamma\le C_Q(N),
 \label{n3:apriori:L2}
\end{equation}
where $\pi$ is projection to the base.  To check this without an
upper bound for $t$, observe that
\[
 e^{2Z}=e^{2t}D^{1/4},\qquad u=le^{2t}\sqrt D,
 \qquad D^{1/4}\le C(N)(1+\tau a\sigma).
\]
The last inequality follows by splitting $l\sigma\le1$ and
$l\sigma>1$, using $\ell\le l\le e$; on the second set
$a\ge1/\sqrt2$ and the right side grows linearly in $\sigma$,
whereas $D^{1/4}$ grows at most as $C\sqrt\sigma$.
Thus the high-level part of \eqref{n3:apriori:L2} follows from
\eqref{n3:apriori:high-energy} and the positive minimum of $\rho$ on
$Q$.  On $Z\le M_0+1$, both $\sigma$ and $e^{2Z}$ are bounded
by the floor identity, which controls the remaining compact integral.

\subsection{Sobolev inequality and iteration on the ordinary graph}

The high-level estimate gives integral control, but degree requires a
pointwise upper bound for $Z$. The graph Sobolev inequality and
iteration supply that bound. The earlier polynomial trace estimate,
not the constants in this iteration, will control the mass limit.

The comparisons in this subsection are uniform in the solution,
$R$, and the homotopy parameter, but may depend arbitrarily on $N$.
If $d_0=(1+\sigma^2)^{-1}$, then
\[
 |d\varphi|_{A_{d_0}}=|\nabla_\Gamma\varphi|,
 \qquad
 \frac{d}{d_0}=\frac{1+\sigma^2}{1+l^2\sigma^2}.
\]
The last ratio is bounded above and below by constants depending
only on $\ell$; Equation~\eqref{n3:apriori:parameter-comparisons}
therefore compares $A_{d_0},A_d,A_\chi$ without any bound for $Z$.

\begin{lemma}[Local graph Sobolev inequality]
\label{n3:apriori:sobolev}
For a smooth function $\omega$ supported over a fixed compact base
patch, which may meet an inner face,
\begin{equation}
 \|\omega\|_{L^6(\dd\Gamma)}^2
 \le C(N)\int_\Gamma
       \left[|\nabla_\Gamma\omega|^2+
                            (1+\mathcal T)\omega^2\right]\dd\Gamma.
 \label{n3:apriori:sobolev-equation}
\end{equation}
\end{lemma}

\begin{proof}
Extend the fixed base patch smoothly and embed a compact enlargement
isometrically in Euclidean space by Nash's embedding theorem
\cite[Theorem~2, p.~59]{Nash1956}.  Its product with the vertical
line gives an isometric Euclidean immersion of the graph.  The
second fundamental form of the base embedding is bounded, so its
contribution to the graph mean curvature vector is bounded.
The scalar mean curvature in the product is
\[
 H_\Gamma=
 \frac{\sqrt D}{l\sqrt{1+\sigma^2}}
       \left(\tr_{e_f^\perp}H^f+d_0H^f(e_f,e_f)\right).
\]
The prefactor is bounded by $C(N)$, and
$d_0/\sqrt d=\sqrt{1+l^2\sigma^2}/(1+\sigma^2)\le e$.
Coercivity therefore gives
$|\boldsymbol H_\Gamma|\le C(N)(1+\sqrt{\mathcal T})$.
This estimate involves controlled quadratic energy, not a pointwise
bound for the Hessian.

The ordinary Michael--Simon inequality \cite{MichaelSimon1973}, in
the compact-support form stated in
\cite[Chapter~4, Section~5, Theorem~5.7, p.~98]{Simon2014}, gives after the
boundary extension described below, for nonnegative $\varphi$,
\[
 \left(\int_\Gamma\varphi^{3/2}\dd\Gamma\right)^{2/3}
 \le C\left[
       \int_\Gamma(|\nabla_\Gamma\varphi|
                          +|\boldsymbol H_\Gamma|\varphi)\dd\Gamma
       +\int_{\partial\Gamma}\varphi\dd A\right].
\]
One can obtain the boundary form directly from the compact-support
version by smoothly extending each smooth graph across its boundary
and cutting off on a collar tending to zero: the cutoff derivative
integral converges to the boundary integral, and the added curvature
integral tends to zero.  This argument is applied to each smooth
graph before the uniform estimate is taken.  The only boundary in
the support is an inner face, where $f$ is constant and
$\dd A=\dd A_g$.  Equation~\eqref{n3:apriori:unweighted-trace} and
the fixed-$N$ graph comparisons bound its contribution by
\[
 C(N)\int_\Gamma
       [(1+\sqrt{\mathcal T})\varphi
                               +|\nabla_\Gamma\varphi|]\dd\Gamma.
\]
Apply the resulting inequality to $\varphi=|\omega|^4$.
Cauchy--Schwarz bounds its right side by
\[
 C(N)\|\omega\|_6^3
 \left(\int_\Gamma
       [|\nabla_\Gamma\omega|^2+(1+\mathcal T)\omega^2]
                                      \dd\Gamma\right)^{1/2}.
\]
The left side is $\|\omega\|_6^4$; division, with the zero case
separate, and squaring prove the assertion.
\end{proof}

To turn the integral bound into an upper bound, let $\eta$ be a
nonnegative cutoff over a fixed compact patch.  For $k$ sufficiently
large, independently of $R$ and the solution, test the first-stage
divergence equation by
\[
 q=\eta^2e^{2kZ}/L.
\]
The positive measure after cancellation of $L$ is
$\sqrt D\dd V_g$.  The differentiated test is
\[
 dq=\frac{e^{2kZ}}L
       [2\eta\,d\eta+\eta^2(2k\,dZ-p_L\,dt)].
\]
Its principal exponential contribution is
$8k\eta^2e^{2kZ}|dZ|_{A_\chi}^2$.  The source is bounded below
by $\delta_0\mathcal T-\Pi_N$ everywhere.
Coercivity and $A_\chi\le A_d$ give
\[
 4p_L|\langle dt,dZ\rangle_{A_\chi}|
 \le\tfrac14\delta_0\mathcal T+C(N)|dZ|_{A_\chi}^2.
\]
Choose $k\ge k_*(N)$ so that the second term is absorbed by the
principal contribution.  The drift terms, with $|K(w,\cdot)|_{A_\chi}
\le |K|$, obey
\begin{align*}
 2k|\langle K(w,\cdot),dZ\rangle_{A_\chi}|
   &\le k|dZ|_{A_\chi}^2+Ck|K|^2,\\
 p_L|\langle K(w,\cdot),dt\rangle_{A_\chi}|
   &\le\tfrac14\delta_0\mathcal T+C(N)|K|^2.
\end{align*}
The cutoff cross terms are handled by the same Cauchy inequalities,
leaving $C(N)k e^{2kZ}(|d\eta|_g^2+\eta^2)$.
The penalty is bounded and contributes to this latter expression.
Finally, Equation~\eqref{n3:apriori:small-boundary-flux} before its last
inequality gives
$|q(V_\lambda)_\nu|\le C\eta^2e^{2kZ}\sqrt d
\le C\eta^2e^{2kZ}$.
Converting to the ordinary graph yields
\begin{align}
 &\int_\Gamma e^{2kZ}\eta^2
       (\mathcal T+k|\nabla_\Gamma Z|^2)\dd\Gamma\notag\\
 &\quad\le C(N)k\int_\Gamma e^{2kZ}
                (\eta^2+|d\eta|_g^2)\dd\Gamma
               +C(N)\int_B e^{2kZ}\eta^2\dd A_g.
 \label{n3:apriori:iteration-energy-boundary}
\end{align}

Use the unweighted trace with $\varphi=e^{2kZ}\eta^2$.  The boundary
term is at most
\[
 C(N)\int_\Gamma e^{2kZ}
   \left[\eta^2(1+\sqrt{\mathcal T}
                   +k|\nabla_\Gamma Z|)+\eta|d\eta|_g\right]
                                                   \dd\Gamma.
\]
Young's inequality absorbs its $\sqrt{\mathcal T}$ term into the
$\mathcal T$ integral and its $k|\nabla_\Gamma Z|$ term into the
$k|\nabla_\Gamma Z|^2$ integral.  The latter absorption costs
$C(N)k\eta^2e^{2kZ}$, which has the allowed size.  Hence
\begin{equation}
 \int_\Gamma e^{2kZ}\eta^2
       (\mathcal T+k|\nabla_\Gamma Z|^2)\dd\Gamma
 \le C(N)k\int_\Gamma e^{2kZ}
                   (\eta^2+|d\eta|_g^2)\dd\Gamma,
 \label{n3:apriori:iteration-energy}
\end{equation}
with constants independent of $k\ge k_*$.

Apply Lemma~\ref{n3:apriori:sobolev} to $\omega=\eta e^{kZ}$ and use
\eqref{n3:apriori:iteration-energy}.  For nested base patches
$U_r\subset U_{r'}$ and a cutoff satisfying
$|d\eta|_g\le C/(r'-r)$, this gives
\begin{equation}
 \|e^Z\|_{L^{6k}(\Gamma|_{U_r})}
 \le\left[C(N)k^2(1+(r'-r)^{-2})\right]^{1/(2k)}
                  \|e^Z\|_{L^{2k}(\Gamma|_{U_{r'}})}.
 \label{n3:apriori:moser-step}
\end{equation}
Here $\Gamma|_U$ denotes the part of the graph projecting to $U$.
Iterate with $k_j=3^jk_*$ and radii decreasing geometrically from
$r'$ to $r$.  The sums $\sum k_j^{-1}$ and $\sum j/k_j$ are finite.
Taking the product in \eqref{n3:apriori:moser-step} therefore gives
\begin{equation}
 S(r):=\sup_{U_r}e^Z
 \le C(N)(r'-r)^{-A(N)}
             \|e^Z\|_{L^{2k_*}(\Gamma|_{U_{r'}})}.
 \label{n3:apriori:moser-high-norm}
\end{equation}
Take all these patches inside a slightly larger fixed patch on
which \eqref{n3:apriori:L2} holds.  Increasing $k_*$ to at least two,
interpolation gives
\[
 \|e^Z\|_{2k_*}
 \le S(r')^{1-1/k_*}
                  \left(\int e^{2Z}\dd\Gamma\right)^{1/(2k_*)}.
\]
Thus $S(r)\le C(N)(r'-r)^{-A(N)}S(r')^{1-1/k_*}$.
For every $\alpha>0$, Young's inequality makes this
\begin{equation}
 S(r)\le\alpha S(r')+
                   C(N,\alpha)(r'-r)^{-A(N)k_*}.
 \label{n3:apriori:moser-interpolation}
\end{equation}
Choose radii increasing geometrically to a fixed outer radius $r_1$,
and then choose $\alpha<2^{-A(N)k_*-1}$.  Iterating
\eqref{n3:apriori:moser-interpolation} has a convergent geometric sum.
Its final remainder $\alpha^jS(r_j)$ tends to zero, because each
individual smooth solution is bounded on the closed outer patch.
The resulting smaller-patch bound is independent of that individual
supremum.  A finite cover proves a uniform compact upper bound for
$Z$, including at every inner face.

Beyond a fixed sphere containing the drift support, the equation is
$\Div(4uA_\chi\nabla Z)=u\Xi$, and $\Xi>0$ above $M_0$.
At an interior maximum above $M_0$ its left side is nonpositive,
because $dZ=0$ and $\Hess Z\le0$.  The outer boundary value is zero.
The compact bound on that fixed sphere consequently extends the
first-stage upper bound to the whole truncation.

\subsection{The remaining stages and the bounded-variable conclusion}

In the second and third stages the drift is zero.  Their source has
the same high-$Z$ positivity as \eqref{n3:apriori:high-source}, with a
uniform threshold for the bounded homotopy range.  By
Lemma~\ref{n3:apriori:collars}, on an inner face
\begin{equation}
 Z=t+\tfrac18\log(1+l^2\sigma^2)
       \le t+\tfrac18\log(1+e^2C_*^2/\tau^2).
 \label{n3:apriori:high-boundary-t}
\end{equation}
Thus sufficiently high boundary $Z$ forces $j(t)=1$.
The prescribed boundary flux is then zero, since $\lambda=0$.
At the face $A_\chi\nu=\chi\nu$, so this is exactly
$\partial_\nu Z=0$.  The maximum principle excludes a high interior
maximum, and the boundary point principle excludes a high maximum
on a smooth inner face with this zero conormal derivative.  Both
principles are applied to the smooth individual solution on its
finite domain; its coefficients are positive definite there.
Locating the maximum requires no uniform ellipticity constant.
The outer value is zero, so these two stages also have a uniform upper
bound.

In the fourth stage $f=0$ and $t=Z$.  When the common source/flux
scale is positive, the same maximum and boundary argument works
with its positive high-level source.  At scale zero the homogeneous
mixed problem and the outer zero Dirichlet value give $Z=0$.
The upper bound is consequently uniform also as the scale tends
to zero.

\begin{proposition}[Bounds before regularity]
\label{n3:apriori:bounds}
Fix the prepared data and $\epsilon>0$.  There are $N_0$ and
$R_{\min}(N)\to\infty$ such that, for each fixed $N\ge N_0$ and
every $R\ge R_{\min}(N)$, all smooth solutions of the four-stage
homotopy with $t\ge-\epsilon$ satisfy
\begin{equation}
 |h|\le C_0,\qquad |f|\le C_0/\tau,\qquad
 -\epsilon\le t\le Z\le C(N),\qquad
 |\nabla f|\le C(N).
 \label{n3:apriori:bounded-variables}
\end{equation}
The constants are uniform in $R$ and in the homotopy parameter.
They include putative floor-contact solutions, which are then
excluded by Proposition~\ref{n3:deformation:floor}.
The fixed signed slopes, the third-stage upper slopes, and the
polynomial trace estimates are those of
Lemmas~\ref{n3:apriori:collars} and~\ref{n3:apriori:trace}.
\end{proposition}

\begin{proof}
It remains to extract the bounds for the conformal variable and the
gradient. Since $Z=t+\log D/8$ and $D\ge1$, the floor gives
$-\epsilon\le t\le Z$.  The upper bound for $Z$ gives
\[
 \sigma^2=\frac{e^{8(Z-t)}-1}{l^2}
       \le\ell^{-2}\bigl(e^{8(C(N)+\epsilon)}-1\bigr).
\]
This proves the asserted fixed-$N$ gradient bound.  The height bounds
are Lemma~\ref{n3:deformation:height}.  The finite collection of
sublevel and collar estimates and the polynomial absorption choose
$N_0$ once.  Enlarge $R_{\min}(N)$ to include the outer estimates and
all fixed compact patches used above. The next section derives Hessian
and continuity estimates from these bounds and uses them to prove
existence.
\end{proof}
